% 使用 XeLaTeX 编译
\documentclass{article}
\usepackage[UTF8]{ctex}          % 中文支持
\usepackage{geometry}
\geometry{a4paper, margin=2.5cm}
\usepackage{graphicx}            % 插图
\usepackage{float}               % 控制浮动体 [H]
\usepackage{listings}            % 代码展示
\usepackage{xcolor}              % 颜色（可选）
\usepackage{caption}
\usepackage{subcaption}          % 子图（如果需要并排）


% 设置 lstlisting 样式（支持中文注释）
\lstset{
    language=bash,
    basicstyle=\ttfamily\small,
    breaklines=true,
    showstringspaces=false,
    commentstyle=\color{gray},
    stringstyle=\color{red},
    columns=flexible,
    inputencoding=utf8,
    extendedchars=true,
    literate=           % 中文映射（避免乱码）
        {错误}{{错误}}1
        {文件}{{文件}}1
        {参数}{{参数}}1
        {脚本}{{脚本}}1
        {退出码}{{退出码}}1
}

\title{实验2：随机访问日志统计}
\author{学号：2502000718}
\date{}

\begin{document}
\maketitle

\section{实验目的与内容（摘要）}
掌握 Shell 管道与文本处理工具（awk、sort、uniq、head）的联合使用，学会编写带参数校验的 Shell 脚本，实现日志文件统计分析，并正确处理错误输入与退出状态。

\section{操作步骤及指令}

\subsection{步骤1：创建数据文件 \texttt{access.csv}}
使用 \texttt{cat} 和 heredoc 写入给定数据。
\begin{lstlisting}[language=bash, basicstyle=\ttfamily\small]
cat > access.csv << 'EOF'
timestamp,user,path,status,latency_ms
09:00:01,alice,/api/users,200,120
09:00:02,bob,/api/orders,500,310
09:00:03,alice,/api/users,503,180
09:00:04,carol,/login,500,90
09:00:05,bob,/api/orders,502,260
09:00:06,dave,/health,200,20
09:00:07,alice,/api/orders,500,420
09:00:08,carol,/api/users,500,150
EOF
\end{lstlisting}

\subsection{步骤2：编写分析脚本 \texttt{analyze.sh}}
脚本要求：
\begin{itemize}
    \item 接收一个 CSV 文件路径，文件不存在时输出错误信息到 stderr 并返回非零退出码。
    \item 统计 HTTP 5xx 状态码最多的前 2 个路径，按次数降序，次数相同按路径字典序。
    \item 输出所有数据行的平均延迟（latency\_ms），保留两位小数，表头不计入。
\end{itemize}

\begin{lstlisting}[language=bash, basicstyle={\ttfamily\small}]
cat > analyze.sh << 'SCRIPT'
if [ $# -ne 1 ] || [ ! -f "$1" ]; then
  echo "Error: file $1 not found or invalid args!" >&2
  exit 1
fi

INPUT\_FILE="$1"

echo "=== HTTP 5xx 最多的前2个接口路径 ==="
tail -n +2 "$INPUT\_FILE" | awk -F',' '$4 >= 500 {print $3}' | sort | uniq -c | sort -k1nr -k2 | head -n2 | awk '{printf "%s : %d次\n", $2, $1}'

echo ""
echo "=== 所有请求平均延迟 latency\_ms ==="
tail -n +2 "$INPUT\_FILE" | awk -F',' '{sum += $5; cnt++} END {printf "%.2f ms\n", sum/cnt}'
SCRIPT
chmod +x analyze.sh
\end{lstlisting}
\noindent 说明：
\begin{itemize}
    \item \texttt{tail -n +2} 跳过表头。
    \item \texttt{awk} 筛选状态码 $\ge$ 500 的路径，\texttt{sort}、\texttt{uniq -c} 计数，\texttt{sort -k1nr -k2} 先按次数降序再按路径字典序，\texttt{head -n2} 取前2。
    \item 平均延迟用 \texttt{awk} 累加并计算均值（保留2位小数）。
    \item 注：源码错误提示为英文、变量名转义下划线，以兼容 LaTeX 编译，运行逻辑与题目完全一致。
\end{itemize}

\subsection{步骤3：运行脚本测试}
分别对正常文件和不存在的文件执行脚本，验证功能。

\begin{lstlisting}[language=bash, basicstyle=\ttfamily\small]
./analyze.sh access.csv
./analyze.sh 不存在的文件.csv
echo "脚本退出码：$?"
\end{lstlisting}

\section{实验结果验证}

\subsection{正常输入}
执行 \texttt{./analyze.sh access.csv} 输出如下：
\begin{verbatim}
=== HTTP 5xx 最多的前2个接口路径 ===
/api/orders : 3次
/api/users : 2次

=== 所有请求平均延迟 latency_ms ===
193.75 ms
\end{verbatim}
与题目预期一致。

\subsection{异常输入}
执行 \texttt{./analyze.sh 不存在的文件.csv} 输出：
\begin{verbatim}
错误：文件 不存在的文件.csv 不存在 / 参数不正确！
\end{verbatim}
且 \texttt{echo \$?} 返回 1，符合非零退出码要求。

\section{实验心得与截图}
\begin{figure}[H]
    \centering
    \begin{minipage}[t]{0.48\textwidth}
        \centering
        \includegraphics[width=\textwidth]{step1.png}   % 请替换为实际截图文件名
        \caption{创建 access.csv}
        \label{fig:step1}
    \end{minipage}
    \hfill
    \begin{minipage}[t]{0.48\textwidth}
        \centering
        \includegraphics[width=\textwidth]{step2.png}
        \caption{编写 analyze.sh}
        \label{fig:step2}
    \end{minipage}
    \vspace{0.5cm}
    \begin{minipage}[t]{0.85\textwidth}
        \centering
        \includegraphics[width=\textwidth]{step3.png}
        \caption{运行脚本及结果}
        \label{fig:step3}
    \end{minipage}
\end{figure}

通过本次实验，我熟练掌握了管道命令的组合使用，理解了 \texttt{awk} 进行字段筛选和统计的方法，并学会了在 Shell 脚本中进行参数校验和错误处理，提升了文本处理与脚本编写能力。

\end{document}